% The Slide Definitions
%document
\documentclass[10pt]{beamer}
%theme
\usetheme{metropolis}
% packages
\usepackage{color}
\usepackage{listings}
\usepackage[ngerman]{babel}
\usepackage[utf8]{inputenc}
\usepackage{multicol}


% color definitions
\definecolor{mygreen}{rgb}{0,0.6,0}
\definecolor{mygray}{rgb}{0.5,0.5,0.5}
\definecolor{mymauve}{rgb}{0.58,0,0.82}

\lstset{
    backgroundcolor=\color{white},
    % choose the background color;
    % you must add \usepackage{color} or \usepackage{xcolor}
    basicstyle=\footnotesize\ttfamily,
    % the size of the fonts that are used for the code
    breakatwhitespace=false,
    % sets if automatic breaks should only happen at whitespace
    breaklines=true,                 % sets automatic line breaking
    captionpos=b,                    % sets the caption-position to bottom
    commentstyle=\color{mygreen},    % comment style
    % deletekeywords={...},
    % if you want to delete keywords from the given language
    extendedchars=true,
    % lets you use non-ASCII characters;
    % for 8-bits encodings only, does not work with UTF-8
    frame=single,                    % adds a frame around the code
    keepspaces=true,
    % keeps spaces in text,
    % useful for keeping indentation of code
    % (possibly needs columns=flexible)
    keywordstyle=\color{blue},       % keyword style
    % morekeywords={*,...},
    % if you want to add more keywords to the set
    numbers=left,
    % where to put the line-numbers; possible values are (none, left, right)
    numbersep=5pt,
    % how far the line-numbers are from the code
    numberstyle=\tiny\color{mygray},
    % the style that is used for the line-numbers
    rulecolor=\color{black},
    % if not set, the frame-color may be changed on line-breaks
    % within not-black text (e.g. comments (green here))
    stepnumber=1,
    % the step between two line-numbers.
    % If it's 1, each line will be numbered
    stringstyle=\color{mymauve},     % string literal style
    tabsize=4,                       % sets default tabsize to 4 spaces
    % show the filename of files included with \lstinputlisting;
    % also try caption instead of title
    language = Python,
	showspaces = false,
	showtabs = false,
	showstringspaces = false,
	escapechar = ,
}

\def\ContinueLineNumber{\lstset{firstnumber=last}}
\def\StartLineAt#1{\lstset{firstnumber=#1}}
\let\numberLineAt\StartLineAt



\newcommand{\codeline}[1]{
	\alert{\texttt{#1}}
}


% Author and Course information
% This Document contains the information about this course.

% Authors of the slides
\author{Claas de Boer, Tilman Hinnerichs}

% Name of the Course
\institute{Python-Grundlagen}

% Fancy Logo 
\titlegraphic{\hfill\includegraphics[height=1.25cm]{../images/fsr_logo_cropped}}

% Custom Bindings
% \newcommand{\codeline}[1]{
%	\alert{\texttt{#1}}
%}


% Presentation title
\title{Kontrollstrukturen}
\date{12.11.2020}

\begin{document}

\maketitle

\begin{frame}{Gliederung}
    \setbeamertemplate{section in toc}[sections numbered]
    \tableofcontents
\end{frame}

\section{Wie man ein Python Programm schreibt}
\begin{frame}{Die Python Shell}
	Die Python shell
	\begin{itemize}
		\item erlaubt es schnell Code zu testen
		\item hilft schnell Dinge auszurechnen
		\item ist kein Python Programm
	\end{itemize}
\end{frame}

\begin{frame}[fragile]{Erste Schritte}
	\begin{lstlisting}
		>>> print("Hello World!")
		Hello World!
		>>> 3+3
		6
		>>> print(3+3)
		6
		>>> 3 * 3
		9
		>>> 3 / 3 
		1.0
		>>> 3 ** 3
		27
		>>> 4 % 3
		1
	\end{lstlisting}
\end{frame}

\begin{frame}[fragile]{Variablenzuweisungen}
	\begin{lstlisting}
		>>> a = 5
		>>> a
		5
		>>> b = a+2
		>>> b
		7
		>>> a*b
		35	
	\end{lstlisting}
\end{frame}

\begin{frame}{DIY: Python Programm schreiben}
	\begin{enumerate}
		\item Öffnet den Editor eurer Wahl
		\item Schreibt euren Programmcode
		\item Speichert die Datei als *.py ab
		\item Führt die Datei mit dem Interpreter eurer Wahl aus
	\end{enumerate}
	Oder
	\begin{enumerate}
		\item Macht dies alles zusammen in der Entwicklungsumgebung eurer Wahl, z.B.: \url{https://repl.it/languages/Python3}
	\end{enumerate}
	\pause Demonstriere Shell vs. Script
\end{frame}

\begin{frame}
	\huge Fragen?
\end{frame}

\begin{frame}{Einige Aufgaben}
	Schreibe ein Python Programm, welches
	\begin{enumerate}
		\item die Fläche eines Kreises mit $r=4cm$ und $\pi = 3.14$ berechnet.
		\item die Zahlen von 1 bis 20 addiert.
	\end{enumerate}
\end{frame}

% ---- Einführung Kontrollstrukturen ----
\section{Was sind Kontrollstrukturen?}
\begin{frame}{Kontrollstrukturen: Wozu?}
    \begin{itemize}
        \item Strukturieren Programmverhalten
        \item für alle imperativen Sprachen ähnlich
        \item z.B.: Verzweigung, Zählschleife
    \end{itemize}
\end{frame}



\begin{frame}{Kontrollstrukturen in Python}
    Python kennt:
    \begin{itemize}
        \item Ein- / Zweiseitige Verzweigung: \texttt{if .. : / if .. : .. else: ..}
        \item Zählschleife: \texttt{for .. in .. :}
        \item Kopfgesteuerte Schleife: \texttt{while .. :}
    \end{itemize}
	\pause
	Darüber hinaus heute:
	\begin{itemize}
		\item Kommentare
		\item Komparatoren ($<,>,==, <=, >=$)
	\end{itemize}
\end{frame}
% ----------------------- Der Python Interpreter ------------------------------
\section{Verzweigungen ("{}If-Statements"{})}
\begin{frame}[fragile]{Einseitige Verzweigungen}
	\begin{lstlisting}
		>>> 5==5
		True
		>>> 5>6 and 5==5
		False
	\end{lstlisting}
	\pause
    \begin{lstlisting}
    if a == b:
        print ("a is equal to b!")
    print ("this is printed every time. bye.")
    \end{lstlisting}
	\pause
	\textbf{Beachte Indentation/Einrückung}
	\begin{lstlisting}
		if a == b:
			print ("a is equal to b!")
			print ("this is not printed every time. bye.")
	\end{lstlisting}
\end{frame}

\begin{frame}[fragile]{Zweiseitige Verzweigungen}
    \begin{lstlisting}
    if a == b:
        print ("a is equal to b!")
    else:
        print ("a is not equal to b!")
        # this is a comment
        # be careful: 1 == 1, but 1 != "1" !
    
    print ("this is printed every time. bye.")
    \end{lstlisting}
\end{frame}

\begin{frame}[fragile]{Verschachtelte Verzweigungen}
    \begin{lstlisting}
    if a < b:
        print ("a is less than b!")
    else:
        if a == b:
            print ("a == b!")
        print ("a >= b!")
        # is also executed when a == b! 
    \end{lstlisting}
\end{frame}

\begin{frame}[fragile]{Kurzform elif Verzweigungen}
    Oder mit \texttt{elif}:
    \begin{lstlisting}
    if a < b:
        print ("a is less than b!")
    elif a == b:
        print ("a == b!")
    else:
        print ("a > b!")
        # is not if when a == b! 
    \end{lstlisting}
\end{frame}

\section{While-Schleife}

\begin{frame}[fragile]{While-Schleife}
    $\rightarrow$ Wiederholt, solange Bedingung erfüllt ist:
    \begin{lstlisting}
    a = 10
    while a >= 0:
        a -= 1
    \end{lstlisting}
\end{frame}

\begin{frame}[fragile]{Vorzeitiges Abbrechen} 
    \begin{lstlisting}
    a = 10
    while a > -1:
        a -= 1
        if a > 3:
            continue # jump to next interation
        print ("countdown:", a) # only printed for 3,2,1
        if a == 1:
            break # break out of loop immediately
    \end{lstlisting}
\end{frame}

\begin{frame}{Zusammenfassung}
	\begin{itemize}
		\item Indentation/Einrückung
		\item Kommentare
		\item Vergleichen von Werten
		\item Boolesche Werte (True, False, and, or)
		\item Kontrollstrukturen
		\begin{itemize}
			\item if, else, elif
			\item while und continue
		\end{itemize}
	\end{itemize}
\end{frame}

\begin{frame}{Einige Aufgaben 2}
	Schreibe ein Python Programm, welches
	\begin{enumerate}
		\item für eine Variable $a$ testet, ob diese zwischen 100 und 2000 liegt.
		\item für eine Variable $a$ testet, ob diese gerade ist.
		\item alle gerade Zahlen zwischen 1 und 10 ausgibt.
		\item zwei Variablen addiert. Wenn die Summe zwischen 15 und 20 liegt, gebe 20 zurück.
	\end{enumerate}
	und printe eine passende Nachricht an den User.
\end{frame}

% nothing to do from here on
\end{document}
